\documentclass[11pt,a4paper]{article}
\usepackage[margin=1in]{geometry}
\usepackage{booktabs}
\usepackage{longtable}
\usepackage{array}
\usepackage{hyperref}
\usepackage{xcolor}
\usepackage{amsmath}
\usepackage{enumitem}
\hypersetup{colorlinks=true,linkcolor=blue,urlcolor=blue,citecolor=blue}

\title{Replication Report: Ocejo et al.\ (2021)\\
\large Whole genome-based characterisation of antimicrobial resistance and genetic diversity in \textit{Campylobacter jejuni} and \textit{Campylobacter coli} from ruminants}
\author{Ollie (OpenClaw AI) --- BVBRC Replication Project (target BVBRC-52)}
\date{Report Date: 2026-07-02}

\begin{document}
\maketitle

\begin{center}
\textbf{Paper:} Ocejo M, Oporto B, Lav\'in JL, Hurtado A. \textit{Scientific Reports} 11:8998 (2021).\\
\textbf{DOI:} \href{https://doi.org/10.1038/s41598-021-88318-0}{10.1038/s41598-021-88318-0} \quad
\textbf{PMID:} 33903652 \quad \textbf{PMC:} PMC8076188\\
Open access: CC BY 4.0 (full text + Supplementary Tables S1--S4 via Europe PMC).\\
\textbf{Compute:} uicgpu (8$\times$A100); envs \texttt{bvbrc14} (abricate/mlst/blast), \texttt{bvbrc38} (SPAdes/fastp), \texttt{bvbrc28} (NCBI datasets).\\
\textbf{LLM-judge:} Argo \texttt{gpt-5.2}; cross-checked with \texttt{claude-opus-4.8}.
\end{center}

\vspace{0.5em}
\noindent\fbox{\parbox{\linewidth}{\textbf{Verdict: PARTIAL REPLICATION (strong).} On a 16-of-70 representative isolate subset, real raw sequencing reads were downloaded from ENA and \emph{de novo} reassembled, then independently genotyped. Assembly statistics match to $\leq 0.8\%$ genome length and $\leq 0.1\%$ GC; MLST sequence types match the paper 16/16 exactly; both headline chromosomal resistance mutations (gyrA Thr86Ile and 23S rRNA A2075G) are reproduced with correct phenotype tracking (including paper's noted CIP-R / gyrA-WT exception C0268); AMR gene content is consistent. Genotype$\leftrightarrow$phenotype concordance is 91.1\% raw / 93.8\% corrected. PARTIAL rather than full REPLICATED because it is a designed 16/70 subset and full-cohort population-structure / phylogeny / statistical-association analyses were not re-run.}}

\section{The Paper's Claims}

The authors selected \textbf{70 \textit{Campylobacter} isolates (40 \textit{C.\ jejuni} + 30 \textit{C.\ coli})} from a ruminant collection (beef cattle, dairy cattle, sheep; northern Spain), with pre-existing phenotypic MICs against six antimicrobials (GEN, STR, TET, CIP, NAL, ERY; plus AMP by E-test). They performed Illumina NovaSeq WGS (2$\times$150 bp), QC (FastQC/Trimmomatic/PRINSEQ), SPAdes v3.13 assembly, QUAST QC, then screened assemblies with BLASTn + ABRicate v1.0 against ResFinder / NCBI / ARG-ANNOT / CARD / MEGARes, called chromosomal point mutations with PointFinder, and typed by 7-gene MLST. Central findings: (i) extensive genetic diversity cleanly separating the two species; (ii) acquired AMR genes and chromosomal mutations whose presence correlated with phenotype; (iii) specific mechanisms --- gyrA Thr86Ile (quinolones), 23S rRNA A2075G (macrolides), tet(O)/mosaic tet (tetracycline), blaOXA-61-like (\(\beta\)-lactams), aminoglycoside-modifying enzymes; (iv) MLST structure dominated by CC-828 (coli) and CC-21 (jejuni).

\begin{longtable}{c p{7.8cm} l c c}
\caption{Claims tested in this replication.}\\
\toprule
\# & Claim & Type & Public? & Tested? \\
\midrule
\endfirsthead
\toprule \# & Claim & Type & Public? & Tested? \\ \midrule
\endhead
C1 & 70 genomes (40 jejuni / 30 coli) deposited under BioProject PRJNA689687, SRR13362733--802. & Data availability & Yes & \checkmark \\
C2 & Draft genomes $\sim$1.61--1.87 Mb, GC $\sim$30.9\% (jejuni 30.4--30.5; coli 31.2--31.5). & Assembly stats & Yes & \checkmark \\
C3 & Species cleanly separate; MLST STs as reported in Table S1. & Genomic/typing & Yes & \checkmark \\
C4 & Fluoroquinolone (CIP/NAL) resistance $\leftarrow$ chromosomal gyrA Thr86Ile (C257T). & Point mutation & Yes & \checkmark \\
C5 & Macrolide (ERY) resistance $\leftarrow$ 23S rRNA A2075G; absent in susceptibles. & Point mutation & Yes & \checkmark \\
C6 & Tetracycline resistance $\leftarrow$ tet(O) + mosaic tet(O/32/O) (jejuni) / tet(O/M/O) (coli). & AMR gene & Yes & \checkmark \\
C7 & blaOXA-61-like family widespread; blaOXA-489 enriched in ST-827 coli. & AMR gene & Yes & \checkmark \\
C8 & Aminoglycoside resistance $\leftarrow$ aph(2$''$)-Ic (GEN), aadE/ant(6)-Ia (STR); sporadic rpsL K43R. & AMR gene + mut & Yes & \checkmark \\
C9 & Very high phenotype$\leftrightarrow$genotype concordance with minor documented discrepancies. & Integrative & Yes & \checkmark \\
C10 & Full-cohort population structure / phylogeny / statistical ST--AMR associations. & Comparative/stats & Partial & $\triangle$ \\
\bottomrule
\end{longtable}

\section{Method (This Replication)}

All steps ran on real, independently downloaded public data; the paper's supplementary tables were used only as end-of-pipeline comparison ground truth.

\subsection{Retrieval and Panel Choice}
Full text + data-availability section fetched as JATS XML from Europe PMC (PMC8076188), identifying BioProject \textbf{PRJNA689687} (runs SRR13362733--SRR13362802). Supplementary PDF (MOESM1) fetched via Europe PMC \texttt{supplementaryFiles}; Tables S1 (isolate metadata + phenotype/MIC + ST/CC) and S2 (per-isolate raw-read + assembly stats) extracted with \texttt{pdftotext -layout} and parsed to JSON. ENA \texttt{filereport} mapped each SRR $\to$ strain alias (C0xxx) $\to$ species $\to$ FASTQ FTP URL. A representative panel of \textbf{16/70} isolates was chosen to exercise every AMR mechanism plus negative controls: ERY-R coli (23S), GEN-R coli (aph), the blaOXA-489 ST-827 coli, TET-R jejuni (tet + mosaic), multiple gyrA-driven CIP/NAL isolates, a fully susceptible control (C0444), a STR-only discordant case (C0430), and \textbf{C0268} (the paper's documented CIP-R / NAL-S / no-gyrA exception).

\subsection{Assembly, AMR Genotyping, Point Mutations, MLST}
Each isolate: ENA paired FASTQ $\to$ \textbf{fastp} (Q25 sliding, min-len 125), \texttt{--reads\_to\_process} capping input to $\sim$150$\times$ (raw $\sim$1125$\times$) $\to$ \textbf{SPAdes \texttt{--isolate}} $\to$ 200 bp contig filter. Assembly stats via Biopython. \textbf{ABRicate v1.0.0} vs ResFinder/NCBI/CARD/ARG-ANNOT/MEGARes/PlasmidFinder (2026-Apr builds). Point mutations were called directly against \textit{C.\ jejuni} NCTC 11168 (GCF\_000009085.1): \texttt{tblastn} for gyrA residue 86 and rpsL residue 43; \texttt{blastn} 23S vs assemblies with the resistance column pinned empirically by aligning ERY-R vs ERY-S isolates (converged on position 2075, matching the paper's \textit{E.\ coli}-numbered A2075G). MLST via \texttt{mlst --scheme campylobacter} (aspA/glnA/gltA/glyA/pgm/tkt/uncA). A 7$\times$16 = 112-call phenotype-vs-genotype table was scored TP/TN/FP/FN, and an LLM judge (Argo \texttt{gpt-5.2}) rendered the coverage/agreement/verdict assessment.

\section{Results vs Paper}

\paragraph{C1 (data availability) --- MATCH.} ENA returns 70 runs: 40 jejuni + 30 coli; each FASTQ $\sim$330--530 MB gz (consistent with $\sim$1125$\times$).

\paragraph{C2 (assembly stats) --- MATCH ($\leq 0.8\%$ length; $\leq 0.1\%$ GC).} All 16 reassembled genomes lie 0.1--0.8\% shorter than the paper (expected from downsampling) with GC identical to $\pm 0.1\%$; the species-specific GC split (jejuni 30.4--30.5\% vs coli 31.2--31.5\%) independently confirms two-species structure. Contig counts run lower because SPAdes \texttt{--isolate} at 150$\times$ + 200 bp filter yields more contiguous assemblies --- a fragmentation difference, not a content difference.

\paragraph{C3 (MLST) --- 16/16 EXACT.} Every independently-derived 7-gene ST matches Table S1: coli in CC-828 (ST-825/827/2097/1055/1595, ST-827/2097 dominant); jejuni in CC-21/42/206 (ST-21 dominant).

\paragraph{C4 (gyrA T86I) --- MATCH incl.\ paper's exception.} Thr86 in C0025, C0430, and C0268 (paper's exception); Ile86 in all 13 CIP/NAL-resistant isolates. C0268 --- CIP-R yet gyrA WT --- is precisely the isolate the paper singles out and is here independently reproduced.

\paragraph{C5 (23S A2075G) --- PERFECT.} G (mutant) in exactly the four ERY-R isolates C0025, C0140, C0541, C0680; A (WT) in all twelve ERY-S. Column-by-column comparison identified 2075 as the \emph{only} differentiating base in the domain-V macrolide loop.

\paragraph{C6 (tetracycline) --- MATCH with reads-level correction.} tet(O) detected in TET-R isolates (C0025, C0551, C0585, C0612, C0642, C0663, C0673, C0882); mosaic tet(O/32/O) detected in \textit{C.\ jejuni} C0437. Three high-coverage multidrug \textit{C.\ coli} (C0140, C0541, C0680) were assembly-negative for tet(O) despite TET-R; direct raw-read BLAST showed tet(O) abundantly present (148 / 155 / 161 read hits) and absent (0 hits) in the TET-susceptible control C0444 --- an assembly artifact of downsampling, not a paper disagreement.

\paragraph{C7 ($\beta$-lactam) --- MATCH.} blaOXA-193 widely; blaOXA-489 in the ST-827 \textit{C.\ coli} (C0444, C0551, C0663), reproducing the paper's ST-827 enrichment.

\paragraph{C8 (aminoglycosides) --- MATCH.} GEN-R \textit{C.\ coli} carry aph(2$''$)-Ic (+ aph(3$'$)-III in two); STR resistance associated with aadE-Cc / ant(6)-Ia; rpsL K43R detected in C0642 (STR-context), matching sporadic rpsL K43R.

\paragraph{C9 (concordance) --- MATCH (91.1\% raw $\to$ 93.8\% corrected).} 112 calls: TP=55, TN=47, FP=5, FN=5. Per-drug (corrected): NAL 100\%, ERY 100\%, GEN 100\%, CIP 94\%, STR 94\%, TET 100\%, AMP 69\%. Residual discordances are exactly the caveats the paper itself raises (C0268 for CIP; blaOXA-presence-alone over-calls AMP because promoter mutation is separately required).

\paragraph{C10 (full-cohort diversity/phylogeny/stats) --- NOT RE-RUN.} Figs 1--3 (core-genome/ST networks, 70-isolate phylogeny, ST$\leftrightarrow$AMR association) were outside subset scope.

\section{Genuine Critique}

This section is intentionally uncharitable. The verdict is PARTIAL, and PARTIAL means genuine gaps.

\begin{enumerate}[leftmargin=1.4em,itemsep=0.35em]
  \item \textbf{Subset scope is a real limitation, not a formality.} 16 of 70 isolates is 23\%. The choice was mechanism-representative --- ensuring every AMR mechanism, the ST-827 blaOXA-489 pocket, and the C0268 exception were exercised --- but it is not a statistical sample. The paper's population-genetics claims (Figs 1--3: species-level diversity, ST-frequency distribution, phylogenetic topology, ST$\leftrightarrow$AMR associations) rest on all 70; nothing in this replication independently supports those. A full-cohort rerun would be the honest next step and is deferred here for compute-time reasons.
  \item \textbf{Downsampling caused a real (though correctable) genotyping failure.} Three \textit{C.\ coli} (C0140, C0541, C0680) were assembly-negative for tet(O) despite being TET-R. Raw-read BLAST rescued the calls (148/155/161 hits versus 0 in the susceptible control), but the assembly step itself lost signal. This is why the raw concordance is 91.1\% instead of 93.8\%. It also means any downstream analysis that relies solely on assembly (which is exactly what ABRicate does, and exactly what the paper's own pipeline did) will inherit the same class of failure at low coverage. Reporting the corrected 93.8\% is honest because reads-level confirmation is a legitimate orthogonal check; but on my \emph{assembly} pipeline alone, the number is 91.1\%.
  \item \textbf{Point mutations were called via a substitute pipeline, not PointFinder.} tblastn/blastn against NCTC 11168 works and produces clean results, but it is not the paper's tool. The internal validation (clean R/S separation; exact position numbering match) is strong evidence that the substitution behaves equivalently on the mutations tested, but I cannot claim tool-level identity. A PointFinder rerun would close this gap.
  \item \textbf{Phenotypes are the paper's, not independent.} No wet-lab MICs were re-generated. The concordance number therefore measures ``my genotype calls vs the paper's phenotype calls'', not ``genotype vs an independent phenotype ground truth.'' Any systematic error in the paper's MICs propagates unaltered into my concordance metric.
  \item \textbf{The AMP 69\% concordance is a real weakness that the concordance framing partially papers over.} The paper's caveat that blaOXA presence alone does not confer AMP resistance (a promoter mutation is required) is real, but it means my rule (gene present $\Rightarrow$ predicted R) is genuinely wrong for AMP. Attributing the 5 FPs to ``documented caveat'' is factually correct but is also a way of saying I did not build the correct AMP-prediction rule. A more careful replication would add the promoter check.
  \item \textbf{The C0268 reproduction is a nice confirmation but not evidence against alternative mechanisms.} C0268 being CIP-R with gyrA WT reproduces the paper. Neither the paper nor I explain \emph{why} --- CmeABC efflux upregulation, unknown novel mutation, phenotypic assay noise, or something else are all in play and untested.
  \item \textbf{No LLM-judge blinding.} The judge saw a summary that already framed the result as strong-partial; verdict-elicitation was not blind. The 2/2 opus-4.8 cross-check reduces but does not eliminate confirmation bias.
  \item \textbf{Numerical reporting rigor.} All quantitative results (16/16 MLST, per-isolate lengths/GC, 148/155/161 reads-level counts, 91.1\%/93.8\% concordance) come from this run's evidence files, but no formal uncertainty (confidence intervals on concordance, per-position coverage on point-mutation calls) is reported. For a 16-isolate/112-call sample the Wilson 95\% CI on 93.8\% is roughly 87.7--97.0\%, which should be quoted if this number is compared to any other study.
\end{enumerate}

\noindent\textbf{Bottom line of the critique:} what was tested was tested cleanly and reproduces the paper. What was \emph{not} tested (full-cohort phylogeny/statistics, PointFinder tool-identity, independent phenotypes, AMP promoter logic) is not tested. PARTIAL is the honest label.

\section{Threats to Validity}
\begin{itemize}[leftmargin=1.4em,itemsep=0.2em]
  \item Subset (16/70), not full cohort.
  \item $\sim$150$\times$ downsampling directly caused three tet(O) assembly dropouts (resolved at read level).
  \item Phenotypes are the paper's; no independent wet-lab MICs.
  \item Point mutations via tblastn/blastn against NCTC 11168, not via PointFinder itself.
\end{itemize}

\section{LLM-Judge Assessment}
Argo \texttt{gpt-5.2} (free): coverage --- the paper's core testable claims (data availability, assembly stats, MLST, both point-mutation mechanisms, gene-content patterns, concordance logic) were tested; full-cohort diversity/phylogeny/statistics and wet-lab phenotyping were not. Agreement --- high (exact MLST, exact macrolide-mutation tracking, reproduced CIP exception, consistent gene content, 93.8\% corrected concordance in the paper's ``high concordance with documented caveats'' regime). Considered verdict: \textbf{PARTIAL} --- ``A real end-to-end rerun on public reads reproduced the paper's core mechanistic and typing claims with high quantitative concordance, which goes beyond a SPOT-CHECK. It's still PARTIAL because only 16/70 isolates were reanalyzed and the full-cohort diversity/phylogeny and statistical association results were not rerun.''

\section*{Verdict}
\textbf{PARTIAL.}\\[0.5em]
\texttt{WAVE\_RESULT set=BVBRC-52 paper=PMID:33903652 verdict=PARTIAL dir=BVBRC-52-Campylobacter-ruminants-AMR-Ocejo2021 one\_line=De-novo reassembled 16/70 ENA isolates of Ocejo 2021; MLST 16/16 exact, gyrA T86I + 23S A2075G reproduced (incl.\ paper's C0268 CIP exception), AMR gene content consistent, genotype-phenotype concordance 93.8\% --- full-cohort diversity/stats not rerun.}

\end{document}
